\section{Related Work}
\label{sec:related_work}

\subsection{Medical Visual Question Answering}

Medical VQA has been studied extensively since the introduction of dedicated benchmarks. Lau et al. \cite{lau2018dataset} published VQA-RAD, a dataset of 3,515 question-answer pairs generated by clinicians over 315 radiology images sourced from MedPix. Questions span multiple categories including presence, modality, abnormality, and organ identification. The dataset distinguishes between closed-ended (yes/no) and open-ended questions, which require different evaluation strategies.

\subsection{Vision-Language Models}

The Transformer architecture \cite{vaswani2017attention} underpins modern VLMs. CLIP \cite{radford2021clip} and ViT \cite{dosovitskiy2021image} established vision encoders trained with contrastive language-image pre-training. LLaVA \cite{liu2023visual} introduced visual instruction tuning, combining a CLIP vision encoder with a language model (Vicuna) through a learned projection layer. LLaVA-1.5 \cite{liu2024improved} improved upon this with an MLP-based vision-language connector and higher-resolution image inputs (336$\times$336). LLaVA-Med \cite{li2023llavamed} adapted this framework specifically for biomedical applications using PubMed image-text pairs.

\subsection{Parameter-Efficient Fine-Tuning}

LoRA \cite{hu2022lora} decomposes weight update matrices into low-rank factors, reducing trainable parameters by orders of magnitude. For a weight matrix $W_0 \in \mathbb{R}^{d \times k}$, LoRA parameterizes the update as:
\begin{equation}
W = W_0 + \Delta W = W_0 + BA
\end{equation}
where $B \in \mathbb{R}^{d \times r}$ and $A \in \mathbb{R}^{r \times k}$, with rank $r \ll \min(d, k)$. QLoRA \cite{dettmers2024qlora} extends this by quantizing $W_0$ to 4-bit NormalFloat (NF4) representation, further reducing memory requirements while maintaining training through full-precision gradients computed for the adapter weights.
